% (c) Nikita Lisitsa, lisyarus@gmail.com, 2024

\documentclass[handout,10pt]{beamer}

\usepackage[T2A]{fontenc}
\usepackage[russian]{babel}
\usepackage{minted}

\usepackage{graphicx}
\graphicspath{ {./images/} }

\usepackage{adjustbox}

\usepackage{color}
\usepackage{soul}

\usepackage{hyperref}

\usetheme{metropolis}

\definecolor{red}{rgb}{1,0,0}
\definecolor{green}{rgb}{0,0.5,0}
\definecolor{blue}{rgb}{0,0,1}
\definecolor{codebg}{RGB}{29,35,49}
\definecolor{lightbg}{RGB}{253,246,227}
\setminted{fontsize=\footnotesize}

\makeatletter
\newcommand{\slideimage}[1]{
  \begin{figure}
    \begin{adjustbox}{width=\textwidth, totalheight=\textheight-2\baselineskip-2\baselineskip,keepaspectratio}
      \includegraphics{#1}
    \end{adjustbox}
  \end{figure}
}
\makeatother

\title{Фотореалистичный рендеринг\quad\quad\quad\quad\quad\quad \textit{(aka raytracing)}}
\subtitle{Лекция 7: Интерполяция нормалей, модели материалов, microfacets}
\date{2024}

\setbeamertemplate{footline}[frame number]

\begin{document}

\frame{\titlepage}

\begin{frame}
\frametitle{Flat/smooth shading}
\begin{itemize}
\item Различают два способа задания нормалей объектов:
\pause
\begin{itemize}
\item \textbf{Flat shading}: нормали вычисляются как честные перпендикуляры к треугольникам
\pause
\item \textbf{Smooth shading}: нормали задаются в вершинах и интерполируются внутри треугольников
\end{itemize}
\end{itemize}
\end{frame}

\begin{frame}
\frametitle{Flat shading}
\slideimage{sphere-flat.png}
\end{frame}

\begin{frame}
\frametitle{Smooth shading}
\slideimage{sphere-smooth.png}
\end{frame}

\begin{frame}
\frametitle{Flat/smooth shading}
\begin{itemize}
\item Особенности \textbf{flat shading}'а:
\pause
\begin{itemize}
\item Отчётливо видны треугольники модели из-за разрывов в освещённости
\pause
\item Не нужно хранить нормали, так как они восстанавливаются по позициям вершин
\end{itemize}
\pause
\item Особенности \textbf{smooth shading'а}:
\pause
\begin{itemize}
\item Объект выглядит гладко, создаётся иллюзия бóльшей детализации
\pause
\item Нужно хранить нормали в вершинах и интерполировать их внутри треугольников
\end{itemize}
\end{itemize}
\end{frame}

\begin{frame}
\frametitle{Flat/smooth shading}
\begin{itemize}
\item На практике всегда используется \textbf{smooth shading}, так как он лучше ложится на GPU:
\pause
\begin{itemize}
\item Нормаль удобно хранить как атрибут вершины
\pause
\item Интерполяцию GPU делает за нас
\end{itemize}
\pause
\item \textbf{Flat shading} на GPU требует дублирования вершин (две вершины с одинаковой позицией, но разными нормалями), либо геометрических шейдеров
\end{itemize}
\end{frame}

\begin{frame}
\frametitle{Flat/smooth shading}
\begin{itemize}
\item В типичной модели могут встретиться как гладкие участки, так и резкие углы
\pause
\item Обычно это достигается дублированием вершин в местах сгибов
\pause
\item В Blender'е этому соответствует режим \textit{auto smooth}
\end{itemize}
\end{frame}

\begin{frame}
\frametitle{Flat/smooth shading в рейтрейсинге}
\begin{itemize}
\item Чтобы поддержать гладкие нормали в рейтрейсинге, нам нужно
\pause
\begin{itemize}
\item 1. Прочитать нормали из входной сцены
\pause
\item 2. Проинтерполировать их внутри треугольников
\pause
\item 3. Правильно учитывать их в расчёте освещения
\end{itemize}
\end{itemize}
\end{frame}

\begin{frame}
\frametitle{Flat/smooth shading в рейтрейсинге}
\usemintedstyle{solarized-light}
\begin{itemize}
\item В формате glTF нормали описаны как атрибут вершин \mintinline{cpp}|NORMAL|, так же, как атрибут позиции вершин \mintinline{cpp}|POSITION|
\pause
\item Так же, как с позициями, будем считать, что это всегда трёхмерный вектор из 32-битных floating point чисел (\mintinline{cpp}|type == "VEC3"|, \mintinline{cpp}|componentType == 5126|)
\pause
\item Можно сохранить нормали в отдельном массиве (параллельно позициям вершин), а можно завести структуру для вершины
\pause
\item Учтите, что в дальнейшем у нас появится ещё один атрибут вершин -- текстурные координаты
\end{itemize}
\end{frame}

\begin{frame}
\frametitle{Flat/smooth shading в рейтрейсинге}
\usemintedstyle{solarized-light}
\begin{itemize}
\item Есть два варианта организации интерполяции нормалей:
\pause
\begin{itemize}
\item Возвращать из функции пересечения луча с треугольником параметры линейной интерполяции, т.е. барицентрические координаты (\begin{math}u, v\end{math} из 5ой лекции), и вычислять интерполированную нормаль там, где это нужно
\pause
\item Сразу вычислять интерполированную нормаль в функции пересечения
\end{itemize}
\pause
\item Делайте, как вам нравится, но мне кажется, что второй способ удобнее
\end{itemize}
\end{frame}

\begin{frame}
\frametitle{Flat/smooth shading в рейтрейсинге}
\usemintedstyle{solarized-light}
\begin{itemize}
\item У нас есть нормали \begin{math}N_A,N_B,N_C\end{math} трёх вершин треугольника, и барицентрические координаты \begin{math}u,v\end{math} точки внутри треугольника
\pause
\item Как вычислить проинтерполированную нормаль?
\pause
\item Для сферической интерполяции между двумя точками есть явная формула -- \mintinline{cpp}|slerp| (\textit{Spherical Linear intERPolation})
\pause
\item Для бóльшего числа точек есть неприятные итеративные алгоритмы, не подходящие из-за низкой производительности
\pause
\item Обычно берут линейную интерполяцию нормалей, и нормируют получившийся вектор:
\begin{gather*}
N' = N_A + u\cdot (N_B - N_A) + v \cdot (N_C - N_A) \\
N = \frac{N'}{\|N'\|}
\end{gather*}
\end{itemize}
\end{frame}

\begin{frame}
\frametitle{Flat/smooth shading в рейтрейсинге}
\usemintedstyle{solarized-light}
\begin{itemize}
\item Просто заменить плоскую нормаль на гладкую -- неправильно, так как некоторые её использования ожидают настоящую нормаль к геометрии:
\pause
\begin{itemize}
\item Вычисление, пришёл луч изнутри или снаружи объекта
\pause
\item Вычисление вероятностей для light surface sampling'а
\end{itemize}
\end{itemize}
\end{frame}

\begin{frame}
\frametitle{Неправильное вычисление внутри/снаружи}
\usemintedstyle{solarized-light}
\slideimage{sphere-wrong.png}
\end{frame}

\begin{frame}
\frametitle{Flat/smooth shading в рейтрейсинге}
\usemintedstyle{solarized-light}
\begin{itemize}
\item Лучше всего возвращать из функции пересечения сразу две нормали: плоскую \textit{geometry normal} и проинтерполированную \textit{shading normal}:
\pause
\begin{itemize}
\item \textit{Geometry normal} используется для вычислений, связанных с геометрической формой объекта
\pause
\item \textit{Shading normal} используется для вычислений результирующей освещённости (в формулах BRDF, отражённых/преломлённых векторов, направляющего вектора полусферы направлений, и т.д.)
\end{itemize}
\end{itemize}
\end{frame}

\begin{frame}
\frametitle{Geometry (flat) normals}
\usemintedstyle{solarized-light}
\slideimage{geometry_normals.png}
\end{frame}

\begin{frame}
\frametitle{Shading (smooth) normals}
\usemintedstyle{solarized-light}
\slideimage{shading_normals.png}
\end{frame}

\begin{frame}
\frametitle{Модели материалов}
\usemintedstyle{solarized-light}
\begin{itemize}
\item До этого мы делили материалы на три типа (диффузные, металлы, диэлектрики) и обрабатывали их отдельно
\pause
\item Хочется научиться поддерживать более интересные материалы: шероховатый металл, пластик (непрозрачный диэлектрик), и т.п.
\pause
\item Хочется использовать более реалистичные модели материалов
\pause
\item Хочется по возможности иметь одну модель, умеющую всё это, а не обрабатывать разные случаи
\end{itemize}
\end{frame}

\begin{frame}
\frametitle{Модели материалов}
\usemintedstyle{solarized-light}
\begin{itemize}
\item Вспомним, какие методы нам нужны от описания материала в конкретной точке:
\pause
\begin{itemize}
\item \mintinline{cpp}|float brdf(vec3 L, vec3 V)| -- значение BRDF по входному (L, light) и выходному (V, view) направлениям
\pause
\item \mintinline{cpp}|vec3 sample(vec3 V)| -- генератор случайных семплов направлений для importance sampling'а этого материала
\pause
\item \mintinline{cpp}|float pdf(vec3 L)| -- плотность вероятности соответствующего семпла
\pause
\item \mintinline{cpp}|float eval(vec3 L, vec3 V) = brdf(L, V) / pdf(L)| -- множитель, с которым семплы в итоге войдут в формулу Монте-Карло интегрирования
\end{itemize}
\end{itemize}
\end{frame}

\begin{frame}
\frametitle{Модели материалов}
\usemintedstyle{solarized-light}
\begin{itemize}
\item В функцию \mintinline{cpp}|brdf| часто добавляют косинус угла падения, так как он зависит от тех же параметров и может позволить упростить формулу, -- и мы тоже это сделаем
\pause
\item Функция \mintinline{cpp}|eval| полезна по тем же причинам -- в ней может многое сократиться, что позволит быстрее её вычислять
\end{itemize}
\end{frame}

\begin{frame}
\frametitle{Ламбертова BRDF}
\usemintedstyle{solarized-light}
\begin{itemize}
\item \mintinline{cpp}|brdf(L, V) = color / PI * max(0, dot(L, N))|
\pause
\item \mintinline{cpp}|sample(V) = normalize(random_unit() + N)|
\pause
\item \mintinline{cpp}|pdf(L) = max(0, dot(L, N)) / PI|
\pause
\item \mintinline{cpp}|eval(L, V) = color|
\end{itemize}
\end{frame}

\begin{frame}
\frametitle{Металлическая BRDF}
\usemintedstyle{solarized-light}
\begin{itemize}
\item Для идеально отражающего металла BRDF является не функцией, а обобщённой функцией
\pause
\item Тем не менее, её можно выразить в этом интерфейсе, зная, что вызываться будут только функции \mintinline{cpp}|sample()| и \mintinline{cpp}|eval()|
\pause
\item \mintinline{cpp}|sample(V) = 2 * N * dot(V, N) - V = R|
\pause
\item \mintinline{cpp}|eval(L, V)| = \mintinline{cpp}|color|
\end{itemize}
\end{frame}

\begin{frame}
\frametitle{Металлическая BRDF: roughness}
\usemintedstyle{solarized-light}
\begin{itemize}
\item В реальности металлические объекты часто отражают не идеально, а с небольшим размытием
\pause
\item Этим управляет параметр \textit{roughness (грубость, шероховатость)}:
\pause
\begin{itemize}
\item \mintinline{cpp}|roughness = 0| -- идеально отполированный материал
\pause
\item \mintinline{cpp}|roughness = 1| -- максимально грубый материал
\end{itemize}
\pause
\item \textbf{\alert{N.B.}}: разные модели материалов по-разному интерпретируют значение \mintinline{cpp}|roughness|
\end{itemize}
\end{frame}

\begin{frame}
\frametitle{Металлическая BRDF: roughness}
\usemintedstyle{solarized-light}
\slideimage{roughness.png}
\end{frame}

\begin{frame}
\frametitle{Металлическая BRDF: roughness}
\usemintedstyle{solarized-light}
\begin{itemize}
\item Есть простой способ поддержать параметр roughness в нашем интерфейсе: прибавить к направлению отражённого луча случайный вектор на сфере радиуса roughness, `размывая' его
\pause
\item \mintinline{cpp}|sample(V) = normalize(R + random_unit() * roughness)|
\pause
\item Формально, мы задали алгоритмом семплинга некое распределение направлений, и потом положили \mintinline{cpp}|brdf| равной \mintinline{cpp}|color * pdf|
\end{itemize}
\end{frame}

\begin{frame}
\frametitle{Металлическая BRDF: roughness}
\usemintedstyle{solarized-light}
\begin{itemize}
\item Это рабочая, но не очень хорошая реализация roughness:
\pause
\begin{itemize}
\item Мы не знаем pdf, который пригодится для importance sampling'а (но при желании её можно явно вычислить)
\pause
\item Это чисто эвристическая модель, и реальные материалы ведут себя иначе
\pause
\item При \mintinline{cpp}|roughness| близком к 1 возникают артефакты
\end{itemize}
\end{itemize}
\end{frame}

\begin{frame}
\frametitle{Importance sampling}
\usemintedstyle{solarized-light}
\begin{itemize}
\item Описанный подход к идеально гладким металлам плохо ложится на multiple importance sampling
\pause
\item В MIS мы генерируем луч одним из двух способов и вычисляем суммарную плотность вероятности
\pause
\item Если одно из распределений -- дискретное, то у него нет плотности вероятности \begin{math}\Longrightarrow\end{math} нам не посчитать суммарную плотность, и MIS неприменим
\pause
\item Проще всего это исправить, ограничив \mintinline{cpp}|roughness| каким-то значением снизу, например, \mintinline{cpp}|0.03|
\end{itemize}
\end{frame}

\begin{frame}
\frametitle{Specular reflection}
\usemintedstyle{solarized-light}
\begin{itemize}
\item Термином \textit{specular} часто называют любое неметаллическое отражение света от поверхности
\pause
\item Одна из первых предложенных specular brdf -- модель Фонга
\pause
\item \mintinline{cpp}|brdf(L, V) = reflectance * max(0, dot(L, N)) * | \mintinline{cpp}|* pow(max(0, dot(R, L)), shininess)|
\pause
\item У неё есть ряд недостатков:
\pause
\begin{itemize}
\item Она эвристическая, и реальные материалы выглядят не так
\pause
\item Её тяжело нормировать (чтобы объект с \mintinline{cpp}|reflectance = 1| отражал 100\% света)
\pause
\item Она не удовлетворяет Helmholtz reciprocity: \mintinline{cpp}|brdf(L, V) != brdf(V, L)| (без учёта \mintinline{cpp}|dot(L, N)|)
\end{itemize}
\end{itemize}
\end{frame}

\begin{frame}
\frametitle{Phong model}
\usemintedstyle{solarized-light}
\slideimage{phong.png}
\end{frame}

\begin{frame}
\frametitle{Blinn-Phong specular}
\usemintedstyle{solarized-light}
\begin{itemize}
\item Модификация specular составляющей модели Фонга
\pause
\item \mintinline{cpp}|brdf(L, V) = reflectance * max(0, dot(L, N)) * | \mintinline{cpp}|* pow(max(0, dot(H, N)), shininess)|
\item \mintinline{cpp}|H = normalize(L + V)| -- \textit{halfway} вектор
\pause
\begin{itemize}
\item Всё ещё эвристическая, но выглядит лучше, чем модель Фонга
\pause
\item Всё ещё тяжело нормировать
\pause
\item \textbf{Удовлетворяет} Helmholtz reciprocity
\end{itemize}
\end{itemize}
\end{frame}

\begin{frame}
\frametitle{Blinn-Phong model}
\usemintedstyle{solarized-light}
\slideimage{blinn-phong.png}
\end{frame}

\begin{frame}
\frametitle{Microfacets}
\usemintedstyle{solarized-light}
\begin{itemize}
\item Современные \textit{physically-based (PBR)} модели материалов основывают на теории \textit{микрофасетов}
\pause
\item Предположим, что на микро-уровне поверхность состоит из маленьких кусочков, каждый со своей нормалью
\pause
\item Предположим, что сами кусочки имеют какую-нибудь простую BRDF -- Ламбертову, или идеального зеркала
\pause
\item Вычислим BRDF макро-объекта, исходя из этой модели
\end{itemize}
\end{frame}

\begin{frame}
\frametitle{Microfacets}
\usemintedstyle{solarized-light}
\slideimage{microfacets.png}
\end{frame}

\begin{frame}
\frametitle{Oren-Nayar model}
\usemintedstyle{solarized-light}
\begin{itemize}
\item Модель \textit{Орена-Наяра} -- микрофасетная BRDF, предполагающая гауссово распределение нормалей микрофасетов с отклонением \begin{math}\sigma\end{math}, и Ламбертову BRDF микрофасетов
\pause
\item Есть много вариаций модели, отличающихся сложностью и точностью вычислений
\end{itemize}
\end{frame}

\begin{frame}
\frametitle{Oren-Nayar model}
\usemintedstyle{solarized-light}
\begin{itemize}
\item \begin{math}f(L,V) = \frac{\text{color}}{\pi}\cdot \max(0, L\cdot N)\cdot (A + B \max(0, \cos\varphi)\sin \alpha \tan \beta)\end{math}
\item \begin{math}A = 1 - 0.5 \frac{\sigma^2}{\sigma^2 + 0.33} \quad B = 0.45 \frac{\sigma^2}{\sigma^2 + 0.09}\end{math}
\item \begin{math}\alpha = \max(\theta_L,\theta_V) \quad \beta = \min(\theta_L,\theta_V)\end{math}
\item \begin{math}\theta_L = \angle(L, N) \quad \theta_V = \angle(V, N)\end{math}
\item \begin{math}\varphi\end{math} -- угол между проекциями \begin{math}L\end{math} и \begin{math}V\end{math} на плоскость, перпендикулярную \begin{math}N\end{math}
\end{itemize}
\end{frame}

\begin{frame}
\frametitle{Oren-Nayar model}
\usemintedstyle{solarized-light}
\slideimage{oren-nayar.jpg}
\end{frame}

\begin{frame}
\frametitle{Disney (Burley) diffuse BRDF}
\usemintedstyle{solarized-light}
\begin{itemize}
\item Простая модель, подогнанная под известные базы данных материалов
\pause
\item \begin{math}f(L,V) = \frac{\text{color}}{\pi}\cdot F_{Schlick}(\max(0, L\cdot N))\cdot F_{Schlick}(\max(0, V\cdot N))\end{math}
\item \begin{math}F_{Schlick}(X) = 1 + (F_{D90} - 1)(1 - X)^5\end{math}
\item \begin{math}F_{D90} = 0.5 + 2\cdot \text{roughness}\cdot \cos^2\theta_d\end{math}
\item \begin{math}\theta_d = \angle(L, H) = \angle(V, H)\end{math}
\end{itemize}
\end{frame}

\begin{frame}
\frametitle{Disney (Burley) diffuse BRDF}
\usemintedstyle{solarized-light}
\slideimage{burley.png}
\end{frame}

\begin{frame}
\frametitle{Cook-Torrance specular BRDF}
\usemintedstyle{solarized-light}
\begin{itemize}
\item Микрофасетная модель, предполагающая, что микрофасеты -- идеальные зеркала
\pause
\begin{equation*}f(L,V) = \frac{D(H)\cdot G(L,V)\cdot F(L, H)}{4 \cdot (L\cdot N)\cdot(V\cdot N)}\end{equation*}
\pause
\item \begin{math}D(H)\end{math} -- \textit{distribution term}: какой процент нормалей микрофасетов ориентированы в сторону \begin{math}H\end{math}
\pause
\item \begin{math}G(L,V)\end{math} -- \textit{geometry term}: насколько соседние микрофасеты мешают свету прийти от источника к поверхности (\textit{shadowing}) или от поверхности к наблюдателю (\textit{masking})
\pause
\item \begin{math}F(L,H)\end{math} -- \textit{fresnel term}: коэффициент отражения Френеля
\pause
\item Обычно \begin{math}D\end{math} зависит от параметра \mintinline{cpp}|roughness|, а \begin{math}G\end{math} вычисляется на основе выбранной функции \begin{math}D\end{math}
\end{itemize}
\end{frame}

\begin{frame}
\frametitle{Beckmann distribution}
\usemintedstyle{solarized-light}
\begin{equation*}
D(H) = \frac{\exp\left(-\frac{\tan^2\theta_H}{\alpha^2}\right)}{\pi\alpha^2\cos^4\theta_H}
\end{equation*}
\begin{itemize}
\item \begin{math}\theta_H = \angle(H,N)\end{math}
\item \begin{math}\alpha = \sqrt{2\sigma}\end{math}
\end{itemize}
\end{frame}

\begin{frame}
\frametitle{GGX distribution}
\usemintedstyle{solarized-light}
\begin{equation*}
D(H) = \frac{\alpha^2\chi^+(H\cdot N)}{\pi\left((\alpha^2-1)(H\cdot N)^2 + 1\right)^2}
\end{equation*}
\begin{itemize}
\item \begin{math}\chi^+(x) = \mathbf 1_{[0, \infty)}(x)\end{math} -- step-функция Хевисайда
\item \begin{math}\alpha = \sqrt{2\sigma}\end{math}
\end{itemize}
\end{frame}

\begin{frame}
\frametitle{Smith's geometry term}
\usemintedstyle{solarized-light}
\begin{gather*}
G_1(N,X) = \frac{1}{1 + \lambda(A)} \\
G = G_1(N,L) \cdot G_1(N,V) \\
G' = \frac{1}{1 + G_1(N,L) + G_1(N,V)}
\end{gather*}
\begin{itemize}
\item \begin{math}A = \frac{(N\cdot X)\cdot\chi^+(N\cdot X)}{\alpha \sqrt{1 - (N\cdot X)^2}}\end{math}
\item \begin{math}\lambda(A)\end{math} зависит от конкретного распределения нормалей \begin{math}D\end{math}
\item \begin{math}G'\end{math} -- \textit{height-correlated} geometry term: утверждается, что такая формула более точно описывает masking и shadowing
\end{itemize}
\end{frame}

\begin{frame}
\frametitle{Smith's geometry term}
\usemintedstyle{solarized-light}
\begin{gather*}
\lambda_{Beckmann}(A) = \mathbf 1_{[0, 1.6]}(A) \cdot \frac{1 - 1.259 A + 0.396 A^2}{3.535 A + 2.181 A^2} \\
\lambda_{GGX}(A) = \frac{1}{2}\left(\sqrt{1 + \frac{1}{A^2}} - 1\right)
\end{gather*}
\end{frame}

\begin{frame}
\frametitle{Fresnel term}
\usemintedstyle{solarized-light}
\begin{gather*}
F = F_0 + (F_{90} - F_0) \cdot (1 - H\cdot L)^5
\end{gather*}
\pause
\begin{itemize}
\item \begin{math}F_{90}\end{math} -- коэффициент отражения при взгляде параллельно поверхности (\textit{grazing incidence}), обычно равен 1
\pause
\item \begin{math}F_0\end{math} -- коэффициент отражения при взгляде прямо на поверхность (\textit{normal incidence})
\pause
\item Для металлов, \begin{math}F_0\end{math} совпадает с цветом материала
\pause
\item Для диэлектриков, \begin{math}F_0\end{math} зависит от коэффициента преломления и колеблется от 0.02 (вода) до 0.18 (алмаз), обычно берут значение 0.04 (соответствующее коэффициенту преломления 1.5)
\end{itemize}
\end{frame}

\begin{frame}
\frametitle{Specular + diffuse}
\usemintedstyle{solarized-light}
\begin{itemize}
\item Коэффициент Френеля \begin{math}F\end{math} показывает, как много света отразится от поверхности
\pause
\item В Cook-Torrance specular BRDF этот коэффициент уже есть
\pause
\item Можно построить BRDF для specular + diffuse, добавив к specular какую-нибудь диффузную BRDF с коэффициентом \begin{math}1-F\end{math}
\pause
\item Тогда коэффициент Френеля можно интерпретировать как коэффициент интерполяции между диффузной и specular BRDF (последняя -- без учёта множителя \begin{math}F\end{math})
\end{itemize}
\end{frame}

\begin{frame}
\frametitle{Sampling microfacet BRDF}
\usemintedstyle{solarized-light}
\begin{itemize}
\item Чем меньше значение roughness, тем ближе specular-часть материала к идеальному зеркалу
\pause
\item \begin{math}\Longrightarrow\end{math} Cosine-weighted sampling плохо подходит для таких материалов
\pause
\item Обычно importance sampling для модели Кука-Торранса семплит нормали в соответствии с распределением D, и использует их для вычисления отражённого луча
\pause
\item Более точным оказывается семплить нормали с учётом \textit{geometry term}, это называется \textit{visible normal distribution (VNDF)}
\end{itemize}
\end{frame}

\begin{frame}
\frametitle{VNDF}
\usemintedstyle{solarized-light}
\begin{itemize}
\item Алгоритм, основанный на преобразовании и семплинге эллипсоидов
\pause
\item В статье \href{https://jcgt.org/published/0007/04/01/paper.pdf}{\texttt{Heitz 2018}} есть все формулы и код для семплинга нормалей
\pause
\item Эту нормаль нужно использовать, чтобы отразить вектор направления луча и получить семпл направления, из которого приходит свет
\pause
\item Обратите внимание, что алгоритм работает в локальной системе координат, где ось \begin{math}Z = N\end{math} совпадает с shading normal, -- нужно будет преобразовать векторы из глобальной системы координат в локальную и обратно
\pause
\item Алгоритм поддерживает \textit{анизотропию} и использует два значения \begin{math}\alpha_X, \alpha_Y\end{math} -- мы положим их равными \begin{math}\alpha\end{math}
\pause
\item В статье также указаны получающаяся pdf и финальное отношение brdf/pdf -- нам нужно только первое, для MIS
\end{itemize}
\end{frame}

\begin{frame}
\frametitle{MIS}
\usemintedstyle{solarized-light}
\begin{itemize}
\item Cosine-weighted хорошо работает для диффузных материалов
\pause
\item VNDF хорошо работает для specular материалов
\pause
\item Light sampling хорошо работает для входящего количества света
\pause
\item \begin{math}\Longrightarrow\end{math} Имеет смысл использовать комбинацию всех трёх с какими-то весами (например, с равными)
\end{itemize}
\end{frame}

\begin{frame}
\frametitle{glTF standard BRDF}
\usemintedstyle{solarized-light}
\begin{itemize}
\item В формате glTF (Appendix B) стандартный материал описывается как смесь диэлектрика и металла, в зависимости от параметра \mintinline{cpp}|metallicFactor|
\pause
\item В качестве specular BRDF берётся Cook-Torrance BRDF с GGX распределением нормалей и uncorrelated Smith's geometry term, но без коэффициента Френеля (по причинам, описанным на предыдущем слайде)
\pause
\item Параметр \begin{math}\alpha\end{math} для GGX полагается равным \begin{math}\text{roughness}^2\end{math}
\pause
\item В качестве диффузной BRDF берётся обычная Ламбертова BRDF: \begin{math}\frac{\text{baseColorFactor}}{\pi}\end{math}
\end{itemize}
\end{frame}

\begin{frame}
\frametitle{glTF standard BRDF}
\usemintedstyle{solarized-light}
\begin{itemize}
\item В качестве металлической BRDF берётся specular BRDF, умноженная на коэффициент Френеля с \begin{math}F_{90}=1\end{math} и \begin{math}F_0=\text{baseColorFactor}\end{math}
\pause
\item В качестве диэлектрической BRDF берётся линейная интерполяция между диффузной и specular BRDF с коэффициентом, равным коэффициенту Френеля для \begin{math}F_{90}=1\end{math} и \begin{math}F_0=0.04\end{math}
\end{itemize}
\end{frame}

\begin{frame}
\frametitle{Сцена без PBR, диффузные материалы}
\usemintedstyle{solarized-light}
\slideimage{pbr-base.png}
\begin{center}1024 spp\end{center}
\end{frame}

\begin{frame}
\frametitle{PBR, specular}
\usemintedstyle{solarized-light}
\slideimage{pbr-specular-50-vndf.png}
\begin{center}1024 spp, roughness = 0.5\end{center}
\end{frame}

\begin{frame}
\frametitle{PBR, specular, без VNDF}
\usemintedstyle{solarized-light}
\slideimage{pbr-specular-0.png}
\begin{center}1024 spp, roughness = 0\end{center}
\end{frame}

\begin{frame}
\frametitle{PBR, specular, c VNDF}
\usemintedstyle{solarized-light}
\slideimage{pbr-specular-0-vndf.png}
\begin{center}1024 spp, roughness = 0\end{center}
\end{frame}

\begin{frame}
\frametitle{PBR, metallic, без VNDF}
\usemintedstyle{solarized-light}
\slideimage{pbr-metallic-50.png}
\begin{center}1024 spp, roughness = 0.5\end{center}
\end{frame}

\begin{frame}
\frametitle{PBR, metallic, с VNDF}
\usemintedstyle{solarized-light}
\slideimage{pbr-metallic-50-vndf.png}
\begin{center}1024 spp, roughness = 0.5\end{center}
\end{frame}

\begin{frame}
\frametitle{PBR, specular + metallic, с VNDF}
\usemintedstyle{solarized-light}
\slideimage{pbr-combined-vndf.png}
\begin{center}1024 spp\end{center}
\end{frame}

\begin{frame}
\frametitle{Расширения glTF BRDF}
\usemintedstyle{solarized-light}
\begin{itemize}
\item Стандартный материал glTF многого не умеет
\pause
\item Не позволяет задать отличный от 1.5 коэффициент преломления (влияет на \begin{math}F_0\end{math} диэлектриков) -- это решается расширением \texttt{KHR\_materials\_ior}
\pause
\item Не позволяет выключить или настроить цвет specular составляющей (например, сделать чисто Ламбертову поверхность) -- это решается расширением \texttt{KHR\_materials\_specular}
\pause
\item Не позволяет описать полупрозрачные материалы и преломление света -- это решается расширением \texttt{KHR\_materials\_transmission}
\pause
\item Мы не будем поддерживать эти расширения из-за временны́х ограничений
\end{itemize}
\end{frame}

\begin{frame}
\frametitle{PBR материалы: ссылки}
\usemintedstyle{solarized-light}
\begin{itemize}
\item \href{https://boksajak.github.io/files/CrashCourseBRDF.pdf}{\texttt{Crash Course in BRDF Implementation}}
\item \href{https://google.github.io/filament/Filament.html}{\texttt{Physically Based Rendering in Filament}}
\item \href{https://media.disneyanimation.com/uploads/production/publication_asset/48/asset/s2012_pbs_disney_brdf_notes_v3.pdf}{\texttt{Physically Based Shading at Disney}}
\item \href{https://seblagarde.files.wordpress.com/2015/07/course_notes_moving_frostbite_to_pbr_v32.pdf}{\texttt{Moving Frostbite to Physically Based Rendering}}
\item \href{https://graphics.pixar.com/library/PhysicallyBasedLighting/paper.pdf}{\texttt{Physically Based Lighting at Pixar}}
\end{itemize}
\end{frame}

\end{document}